\documentclass[11pt,a4paper]{article}

\usepackage[utf8]{inputenc}
\usepackage[T1]{fontenc}
\usepackage{microtype}
\usepackage[margin=2.5cm]{geometry}
\usepackage{amsmath,amssymb}
\usepackage{booktabs}
\usepackage[colorlinks=true,linkcolor=black,citecolor=black,urlcolor=black]{hyperref}
\usepackage{cleveref}

\usepackage[backend=biber,style=numeric-comp,sorting=none]{biblatex}
\addbibresource{../../bibliography/clifford-seminar-group.bib}

\newcommand{\R}{\mathbb{R}}
\newcommand{\Q}{\mathbb{Q}}
\newcommand{\ev}{\operatorname{ev}}
\newcommand{\rad}{\operatorname{rad}}
\newcommand{\Hom}{\operatorname{Hom}}

\title{\Large\bfseries Thesenpapier: Dual Spaces, Multilinear Forms, and
  Quadratic Spaces\\[0.3em]
\large Foundational Session~1 --- Clifford Seminar}
\author{Clifford Seminar Group\\
\small Julian L. Avila-Martinez\\
\small Universidad Distrital Francisco Jos{\'e} de Caldas}
\date{}

\begin{document}

\maketitle

\section{Foundational Definitions}
\label{sec:definitions}

\subsection{Groups, Linear Spaces, and Algebras}
\label{subsec:linear_spaces}

Let $G$ be a set equipped with an internal binary operation
$\cdot : G \times G \to G$.
The pair $(G, \cdot)$ constitutes a group if the operation satisfies the
following three conditions:
\begin{enumerate}
  \item[$\mathcal{G}_1$] Associativity: $(a \cdot b) \cdot c = a \cdot (b \cdot
    c)$ for all $a, b, c \in G$.
  \item[$\mathcal{G}_2$] Identity element: there exists an element $e \in G$
    such that $e \cdot a = a \cdot e = a$ for every $a \in G$.
  \item[$\mathcal{G}_3$] Inverse element: for each $a \in G$, there exists an
    element $a^{-1} \in G$ such that $a \cdot a^{-1} = a^{-1} \cdot a = e$.
\end{enumerate}
If $a \cdot b = b \cdot a$ holds for all $a, b \in G$, the group is abelian.
An abelian group is written additively as $(V, +)$, with identity element $0$
and inverse $-v$.

Let $K$ be a field.
A linear space over $K$ is an additive abelian group $(V, +)$ equipped with an
external scalar multiplication $\cdot : K \times V \to V$, denoted $(\lambda, v)
\mapsto \lambda v$, satisfying the following conditions for all $\lambda, \mu
\in K$ and all $u, v \in V$:
\begin{enumerate}
  \item Distributivity over group addition: $\lambda (u + v) = \lambda u +
    \lambda v$.
  \item Distributivity over field addition: $(\lambda + \mu) v = \lambda v + \mu
    v$.
  \item Compatibility with field multiplication: $\lambda (\mu v) = (\lambda
    \mu) v$.
  \item Unitary action: $1_K v = v$, where $1_K$ denotes the multiplicative
    identity of $K$.
\end{enumerate}
The designation \emph{vector space} is reserved strictly for concrete coordinate
realizations such as $\R^n$ or explicit geometric displacement spaces.

Let $V$ be a linear space over $K$.
The structure $(A, +, \cdot, \bullet)$ is a $K$-algebra if $(A, +, \cdot)$ is a
linear space over $K$ equipped with a bilinear binary operation $\bullet : A
\times A \to A$, satisfying for all $\lambda \in K$ and $u, v, w \in A$ the
compatibility conditions
\begin{equation}
  u \bullet (v + w) = u \bullet v + u \bullet w, \quad
  (u + v) \bullet w = u \bullet w + v \bullet w,
\end{equation}
and
\begin{equation}
  \lambda (u \bullet v) = (\lambda u) \bullet v = u \bullet (\lambda v).
\end{equation}
The algebra is associative if $(u \bullet v) \bullet w = u \bullet (v \bullet
w)$, and unital if there exists an element $1_A \in A$ such that $1_A \bullet u
= u \bullet 1_A = u$ for all $u \in A$.

\subsection{Dual Spaces and the Canonical Pairing}
\label{subsec:dual_spaces}

Let $V$ be a linear space over $K$.
The dual space $V^*$ is the linear space of all linear functionals on $V$,
\begin{equation}
  V^* = \Hom_K(V, K).
\end{equation}
The evaluation of a functional $\alpha \in V^*$ on an element $v \in V$ is
expressed by the bilinear dual pairing
\begin{equation}
  \langle \cdot, \cdot \rangle : V^* \times V \to K, \quad
  \langle \alpha, v \rangle = \alpha(v).
\end{equation}
The double dual space is $V^{**} = \Hom_K(V^*, K)$.
The canonical evaluation mapping $\ev : V \to V^{**}$ defined by
\begin{equation}
  \ev_v(\alpha) = \langle \alpha, v \rangle
\end{equation}
is an injective linear homomorphism.
When $\dim V < \infty$, the map $\ev$ is an isomorphism.
Consequently, elements of $V$ act natively as linear functionals on $V^*$
without auxiliary structure.

\subsection{The Free Tensor Algebra}
\label{subsec:tensor_algebra}

The tensor algebra of a linear space $V$ over $K$ is the graded associative
$K$-algebra
\begin{equation}
  T(V) = \bigoplus_{r=0}^\infty T^r(V),
\end{equation}
where $T^0(V) = K$, $T^1(V) = V$, and $T^r(V) = \bigotimes^r V$ for $r \ge 2$.
Multiplication is governed by the associative tensor product $\otimes$.
The structure $T(V)$ satisfies a universal property: every linear map from $V$
into an associative unital $K$-algebra $A$ extends uniquely to a unital algebra
homomorphism from $T(V)$ into $A$.
The tensor algebra imposes no algebraic relations between distinct elements
of $V$.

\subsection{Quadratic Forms and Polarization}
\label{subsec:quadratic_forms}

Let $V$ be a real linear space.
A map $q : V \to \R$ is a quadratic form if it satisfies $q(\alpha v) = \alpha^2
q(v)$ for all $\alpha \in \R$ and $v \in V$, and the associated symmetric
bilinear form $b : V \times V \to \R$ defined by the polarization identity
\begin{equation}
  b(u, v) = \frac{1}{2} \bigl( q(u + v) - q(u) - q(v) \bigr)
\end{equation}
is linear in each argument.
The radical of the quadratic space $(V, q)$ is the subspace
\begin{equation}
  \rad(V) = \{ u \in V \mid b(u, v) = 0 \text{ for all } v \in V \}.
\end{equation}
The space $(V, q)$ is non-degenerate if $\rad(V) = \{0\}$.

A symmetric bilinear form induces a linear mapping $b^\flat : V \to V^*$ given
by $b^\flat(u) = b(u, \cdot)$.
Non-degeneracy of $q$ is equivalent to $b^\flat$ being an isomorphism, whose
inverse is denoted $b^\sharp : V^* \to V$.

By Sylvester's law of inertia, any finite-dimensional real quadratic space
decomposes into an orthogonal direct sum $V = V_+ \oplus V_- \oplus V_0$ such
that $q|_{V_+}$ is positive definite, $q|_{V_-}$ is negative definite, and
$q|_{V_0} = 0$.
In the non-degenerate case, $\dim V_0 = 0$, and the signature is denoted $(p,
m)$, where $p = \dim V_+$ and $m = \dim V_-$.
The resulting pseudo-Euclidean vector space is denoted $\R^{p,m}$.

\section{Physical Foundations: Kinematics versus Measurement}
\label{sec:physics}

In classical Lagrangian mechanics on a configuration manifold $M$, generalized
velocities $\dot{x}$ belong to the tangent space $T_x M$, which is a linear
space of directional derivations.
In contrast, generalized forces and canonical momenta $p_i =
\partial_{\dot{x}^i} L$ reside in the cotangent space $T_x^* M$, which is the
dual linear space.
The pairing $\langle p, \dot{x} \rangle$ represents the mechanical work rate and
requires no metric structure.

The conventional identification of velocities with momenta in Newtonian
mechanics relies on the kinetic metric $b(u, v) = m \sum_i u^i v^i$, which
supplies the isomorphism $b^\flat : T_x M \to T_x^* M$.
Kinetic energy $T = \frac{1}{2} q(\dot{x})$ is not a linear invariant, but the
evaluation of this specific quadratic form.

In relativistic mechanics, the causal structure of spacetime is established
prior to the introduction of coordinates or curved geometry.
Given the real linear space $V$ of spacetime displacements, the choice of a
non-degenerate quadratic form of signature $(1, 3)$ divides $V$ into timelike
($q(v) > 0$), spacelike ($q(v) < 0$), and lightlike ($q(v) = 0$) elements.
Light cones are the algebraic zero-sets of $q$, demonstrating that relativistic
causality is an algebraic property of quadratic spaces.

\section{Derivation Agenda for the Blackboard Exposition}
\label{sec:agenda}

The 40-minute Referat follows four sequential stages:
\begin{enumerate}
  \item Axiomatic synthesis of linear spaces from additive abelian groups and
    external field actions, followed by the coordinate-free definition of $V^*$
    and the canonical isomorphism $V \cong V^{**}$.
  \item Construction of the free tensor algebra $T(V)$ as the maximal
    multilinear structure, demonstrating its infinite-dimensional, relation-free
    nature.
  \item Introduction of the quadratic form $q$, derivation of the symmetric
    bilinear form via the polarization identity, and the construction of the
    musical isomorphisms $b^\flat$ and $b^\sharp$.
  \item Orthogonal decomposition of non-degenerate quadratic spaces, proof of
    signature invariance, and classification of isotropic rays $q(v) = 0$.
\end{enumerate}

\section{Theses for Plenary Discussion}
\label{sec:theses}

\begin{description}
  \item[Thesis I]
    The standard textbook practice of identifying velocities with canonical
    momenta via Cartesian index identification ($p_i = m \dot{x}_i$) obscures
    the duality between $V$ and $V^*$.
    This obscures the geometric foundation of Hamiltonian mechanics and creates
    conceptual barriers when transitioning to curved manifolds or relativistic
    field theories.

  \item[Thesis II]
    The free tensor algebra $T(V)$ contains excess degrees of freedom that lack
    geometric meaning.
    A physical geometry cannot emerge from multilinear algebra until an
    algebraic ideal generated by measurement---specifically the quadratic form
    $q(v)$---is factored out.

  \item[Thesis III]
    A metric is not an intrinsic property of a linear space, but an auxiliary
    constitutive mapping $b^\flat : V \to V^*$ that converts displacements into
    linear measurement functionals.

  \item[Thesis IV]
    Causal structure in relativistic physics is an algebraic consequence of the
    signature $(1, 3)$ on a four-dimensional linear space, independent of
    differential geometry or pseudo-Riemannian metric tensors.
\end{description}

\section{Preparatory Exercises}
\label{sec:exercises}

\begin{enumerate}
  \item Consider the set $\mathcal{U}$ of coherent physical units generated by
    the base units of the International System of Units (SI).
    Let group addition on $\mathcal{U}$ be defined by physical unit
    multiplication (for example, $\mathrm{m} \cdot \mathrm{s} = \mathrm{m\,s}$
    and $\mathrm{m} \cdot \mathrm{m} = \mathrm{m}^2$), and define external
    scalar multiplication by $\Q$ as unit exponentiation,
    $(\lambda, u) \mapsto u^\lambda$.
    Show that $(\mathcal{U}, \cdot)$ satisfies the axioms of an additive
    abelian group, with the dimensionless unit $1$ serving as identity, and
    verify that this structure satisfies all axioms of a linear space over $\Q$.

  \item Let the ordered set of seven base units $\mathcal{B} = \{\mathrm{m},
      \mathrm{kg}, \mathrm{s}, \mathrm{A}, \mathrm{K}, \mathrm{mol},
    \mathrm{cd}\}$ form a basis for the SI linear space over $\Q$.
    Determine the explicit coordinate tuples in $\Q^7$ corresponding to:
    \begin{enumerate}
      \item The Newton ($\mathrm{N}$), unit of mechanical force.
      \item The Farad ($\mathrm{F}$), unit of electrical capacitance.
      \item The Lumen ($\mathrm{lm}$), unit of luminous flux.
    \end{enumerate}

  \item Let $V$ be a finite-dimensional linear space over $K$, and let
    $\mathcal{E} = \{e_i\}_{i=1}^n$ be a basis for $V$.
    The dual basis $\mathcal{E}^* = \{\omega^j\}_{j=1}^n \subset V^*$ is
    defined by $\langle \omega^j, e_i \rangle = \delta^j_i$.
    Show that the mapping $e_i \mapsto \omega^i$ depends explicitly on the
    choice of basis, confirming that no canonical isomorphism between $V$ and
    $V^*$ exists.
    In contrast, prove that the evaluation map $\ev : V \to V^{**}$ defined by
    $\ev_v(\alpha) = \langle \alpha, v \rangle$ is completely basis-independent
    and canonical.

  \item In classical mechanics, consider a free particle on $\R^3$ described by
    the Lagrangian $L(v) = \frac{1}{2} m b(v, v)$, where $b$ is a symmetric,
    positive-definite bilinear form.
    Show that the definition of canonical momentum $p = \partial_v L$ is
    identical to the action of the musical isomorphism $b^\flat : V \to V^*$.
    Compute $b^\sharp : V^* \to V$ and show that the classical Hamiltonian
    $H(p) = \langle p, v \rangle - L(v)$ corresponds to the inverse quadratic
    form on the dual space.

  \item Consider the real linear space $V = \R^4$ equipped with the Minkowski
    quadratic form $q(x) = (x^0)^2 - (x^1)^2 - (x^2)^2 - (x^3)^2$.
    \begin{enumerate}
      \item Apply the polarization identity to determine the explicit
        symmetric bilinear form $b(x, y)$.
      \item Show that $\rad(V) = \{0\}$, establishing non-degeneracy.
      \item A linear subspace $W \subset V$ is totally isotropic if $q|_W = 0$.
        Prove that the maximum dimension of any totally isotropic subspace in
        $\R^{1,3}$ is $1$, and explain why physical light rays cannot form a
        two-dimensional planar subspace.
    \end{enumerate}
\end{enumerate}

\section{Preparatory Literature}
\label{sec:literature}

Participants are expected to study the following references prior to the
session:
\begin{itemize}
  \item Abstract groups and linear spaces: Fraleigh
    \cite[Chaps.~1.1, 1.2, and~7.33]{fraleighFirstCourseAbstract2021}; Axler
    \cite[Chaps.~1 and~2]{axlerLinearAlgebraDone2015}.
  \item Dual spaces, multilinear forms, and quadratic spaces: Vaz and da
    Rocha \cite[Chap.~1]{vazIntroductionCliffordAlgebras2016}.
\end{itemize}

\printbibliography

\end{document}
